
\begin{center}
    \textbf{ĐẠI HỌC KINH TẾ - LUẬT}\\
    \textbf{KHOA: KINH TẾ ĐỐI NGOẠI}\\
    \vspace{0.5cm}
    \large \textbf{ĐỀ THI HỌC PHẦN: LẬP VÀ THẨM ĐỊNH DỰ ÁN ĐẦU TƯ}\\
    \textit{Thời gian làm bài: 90 phút}
\end{center}

\rule{\textwidth}{1pt}

\section{Phần 1: Chu kỳ dự án và Ý tưởng đầu tư}

\begin{enumerate}[label=\textbf{Câu \arabic*:}, leftmargin=*]

\item Một dự án đang ở giai đoạn "Chạy thử và Nghiệm thu" nhưng phát hiện hệ thống xử lý nước thải không đạt chuẩn. Theo logic chu kỳ dự án, sai sót này phản ánh yếu kém của bước nào?
\begin{enumerate}[label=\Alph*.]
    \item Nghiên cứu phát hiện cơ hội đầu tư.
    \item Thiết kế và lập dự toán thi công.
    \item Thẩm định dự án (Đánh giá \& Quyết định).
    \item Sàng lọc ý tưởng dự án.
\end{enumerate}

\item Doanh nghiệp A phát hiện mỏ khoáng sản mới dựa trên "Tài nguyên thiên nhiên quốc gia" nhưng dự án bị đình trệ do không phù hợp với "Chủ trương nhà nước". Lỗi này thuộc về:
\begin{enumerate}[label=\Alph*.]
    \item Chiến lược phát triển sản xuất kinh doanh của cơ sở.
    \item Bối cảnh Quốc gia.
    \item Tình hình cung thị trường.
    \item Khả năng chuyển hóa thành hoạt động đầu tư.
\end{enumerate}

\item Tại bước "Sản xuất, kinh doanh", một ý tưởng dự án mới nảy sinh. Nhà quản lý cần làm gì đầu tiên để tránh lãng phí?
\begin{enumerate}[label=\Alph*.]
    \item Lập ngay báo cáo nghiên cứu khả thi (FS).
    \item Đối chiếu ý tưởng với nguồn lực và năng lực của chủ thể đầu tư.
    \item Thuê tư vấn thiết kế bản vẽ thi công.
    \item Tiến hành vay vốn ngân hàng.
\end{enumerate}

\item Phát biểu nào sau đây là ĐÚNG về mặt chi phí cơ hội giữa giai đoạn Chuẩn bị và Thực hiện?
\begin{enumerate}[label=\Alph*.]
    \item Chậm trễ ở giai đoạn Chuẩn bị gây thiệt hại lớn hơn vì chiếm 99,5\% vốn.
    \item Chậm trễ ở giai đoạn Thực hiện gây thiệt hại lớn hơn vì vốn đã đổ vào công trình và là "vốn không sinh lời".
    \item Cả hai giai đoạn có mức độ thiệt hại như nhau.
    \item Chậm trễ ở giai đoạn Chuẩn bị không tốn chi phí.
\end{enumerate}

\item Rủi ro lớn nhất khi bỏ qua "Nghiên cứu tiền khả thi" để thực hiện thẳng "Nghiên cứu khả thi" là:
\begin{enumerate}[label=\Alph*.]
    \item Không thể tính được chỉ số IRR.
    \item Thiếu sự sàng lọc các phương án lựa chọn, dẫn đến chọn phương án không tối ưu.
    \item Bị cơ quan nhà nước phạt hành chính.
    \item Không tuyển dụng được nhân sự.
\end{enumerate}

\item Nếu một dự án dừng lại ngay sau bước "Nghiên cứu tiền khả thi", khoản chi phí nào coi là mất trắng?
\begin{enumerate}[label=\Alph*.]
    \item Chi phí thi công xây lắp.
    \item Chi phí lập báo cáo nghiên cứu sơ bộ.
    \item Chi phí mua máy móc thiết bị.
    \item Chi phí trả nợ lãi vay dài hạn.
\end{enumerate}

\item Khi xác định "Cung thị trường", rủi ro phân tích nào thường dẫn đến thất bại khi vận hành?
\begin{enumerate}[label=\Alph*.]
    \item Không tính đến giá trị thanh lý.
    \item Nhầm lẫn giữa nhu cầu tiềm năng và nhu cầu có khả năng thanh toán.
    \item Quên tính chi phí lập dự toán thi công.
    \item Không phân tích bối cảnh tài nguyên.
\end{enumerate}

\item Dự án đang "Thi công xây lắp" thì chủ đầu tư thay đổi công nghệ. Điều này chứng tỏ khâu nào thực hiện kém?
\begin{enumerate}[label=\Alph*.]
    \item Giai đoạn sản xuất kinh doanh.
    \item Bước Nghiên cứu khả thi (khâu nghiên cứu kỹ thuật).
    \item Khâu thanh lý dự án.
    \item Bước hoàn tất thủ tục triển khai.
\end{enumerate}

\item Ý tưởng dự án đầu tư được coi là "khả thi sơ bộ" khi:
\begin{enumerate}[label=\Alph*.]
    \item Đã có bản vẽ thiết kế thi công chi tiết.
    \item Có triển vọng đem lại hiệu quả và chủ đầu tư có khả năng đưa ra quyết định sơ bộ.
    \item Đã được ngân hàng giải ngân 100\% vốn.
    \item Đã hoàn thành lắp đặt máy móc.
\end{enumerate}

\item Tại sao tiến độ là quan trọng trong giai đoạn chuẩn bị?
\begin{enumerate}[label=\Alph*.]
    \item Để sớm nộp thuế.
    \item Vì đây là giai đoạn vốn không sinh lời, kéo dài sẽ tăng chi phí cơ hội.
    \item Để nhanh chóng thanh lý tài sản cũ.
    \item Để tiết kiệm chi phí văn phòng phẩm.
\end{enumerate}

\item Nếu một dự án có "Đầu ra" thuận lợi nhưng "Đầu vào" không thuận lợi, dự án gặp khó khăn gì?
\begin{enumerate}[label=\Alph*.]
    \item Không bán được hàng cho khách hàng.
    \item Khó khăn trong việc cung ứng nguyên liệu hoặc nhân lực.
    \item Không được Nhà nước cấp phép.
    \item Không thể thanh lý tài sản.
\end{enumerate}

\item "Vòng đời dự án" kết thúc khi nào?
\begin{enumerate}[label=\Alph*.]
    \item Khi trả hết nợ gốc ngân hàng.
    \item Khi hoàn thành việc thanh lý các kết quả đầu tư.
    \item Khi dự án đạt công suất cao nhất.
    \item Khi kết thúc giai đoạn thực hiện đầu tư.
\end{enumerate}

\item Dự án năng lượng tái tạo dựa trên "Chiến lược phát triển KT-XH đất nước" thuộc nhóm căn cứ nào?
\begin{enumerate}[label=\Alph*.]
    \item Bối cảnh Quốc gia.
    \item Chiến lược sản xuất kinh doanh của cơ sở.
    \item Tài nguyên thiên nhiên.
    \item Dự kiến hiệu quả đạt được.
\end{enumerate}

\item Sự khác biệt chi tiết giữa "Nghiên cứu cơ hội" và "Nghiên cứu tiền khả thi" là:
\begin{enumerate}[label=\Alph*.]
    \item Nghiên cứu cơ hội chi tiết hơn.
    \item Nghiên cứu tiền khả thi chi tiết hơn nhằm sàng lọc phương án.
    \item Cả hai đều chi tiết như nhau.
    \item Nghiên cứu cơ hội ở trạng thái động.
\end{enumerate}

\item Tại sao cần theo dõi bước "Sử dụng chưa hết công suất" trong giai đoạn Sản xuất kinh doanh?
\begin{enumerate}[label=\Alph*.]
    \item Để giảm thuế TNDN.
    \item Để tinh chỉnh kỹ thuật và thị trường trước khi đẩy lên mức tối đa.
    \item Để kéo dài thời gian thanh lý.
    \item Để tiết kiệm chi phí khấu hao.
\end{enumerate}

\item Nếu dự án bị "treo" ở bước "Hoàn tất thủ tục triển khai", chi phí chuẩn bị sẽ:
\begin{enumerate}[label=\Alph*.]
    \item Được ngân hàng bồi thường 100\%.
    \item Trở thành chi phí chìm và không thể thu hồi.
    \item Tự động chuyển thành vốn điều lệ.
    \item Giảm dần theo thời gian.
\end{enumerate}

\end{enumerate}

\section{Phần 2: Giai đoạn Chuẩn bị đầu tư và Các cấp độ Nghiên cứu}

\begin{enumerate}[label=\textbf{Câu \arabic*:}, leftmargin=*, resume]

\item Tập đoàn đầu tư 5.000 tỷ đồng. Ngân sách tối đa hợp lý cho lập FS và thẩm định là:
\begin{enumerate}[label=\Alph*.]
    \item 25 tỷ đồng. \item 750 tỷ đồng. \item 50 tỷ đồng. \item 500 tỷ đồng.
\end{enumerate}

\item Điểm khác biệt cốt lõi về "Trạng thái nghiên cứu" giữa Pre-FS và FS là:
\begin{enumerate}[label=\Alph*.]
    \item Tiền khả thi nghiên cứu "Động", Khả thi nghiên cứu "Tĩnh".
    \item Cả hai đều nghiên cứu "Tĩnh".
    \item Tiền khả thi nghiên cứu "Tĩnh", Khả thi nghiên cứu "Động".
    \item Cả hai đều nghiên cứu "Động".
\end{enumerate}

\item Nội dung nào thuộc về "Nghiên cứu Kỹ thuật và Nhân sự" trong báo cáo FS?
\begin{enumerate}[label=\Alph*.]
    \item Phân tích Điểm hòa vốn.
    \item Lựa chọn quy trình công nghệ và Sơ đồ tổ chức quản lý.
    \item Tính toán Chỉ số IRR.
    \item Dự báo Cơ cấu nguồn vốn.
\end{enumerate}

\item Sản phẩm của nghiên cứu tiền khả thi (Pre-FS) phải chứng minh được:
\begin{enumerate}[label=\Alph*.]
    \item Dự án đã thu hồi được vốn.
    \item Cơ hội có triển vọng để quyết định cho phép đầu tư sâu hơn.
    \item Máy móc đã được nhập khẩu.
    \item Đã ký được hợp đồng bao tiêu sản phẩm.
\end{enumerate}

\item Tại sao Báo cáo khả thi (FS) được gọi là "Bước sàng lọc cuối cùng"?
\begin{enumerate}[label=\Alph*.]
    \item Vì sau bước này dự án bị thanh lý.
    \item Vì là căn cứ cuối cùng để khẳng định có đầu tư hay không trước khi chi 85-99,5\% vốn.
    \item Vì không cần nghiên cứu lợi ích KT-XH.
    \item Vì độ chính xác thấp.
\end{enumerate}

\item Hậu quả của việc FS sơ sài trong "Nghiên cứu Kỹ thuật" là:
\begin{enumerate}[label=\Alph*.]
    \item Không tìm được khách hàng.
    \item Dự toán chi phí đầu tư bị sai lệch, gây thiếu vốn khi thi công.
    \item NPV luôn bằng 0.
    \item Không tính được thuế TNDN.
\end{enumerate}

\item Yếu tố trọng tâm để tính toán NCF trong nghiên cứu tài chính của FS là:
\begin{enumerate}[label=\Alph*.]
    \item Danh sách các nhà thầu.
    \item Dự kiến dòng thu, dòng chi và phương án huy động vốn.
    \item Bản vẽ mặt bằng nhà xưởng.
    \item Quy hoạch bối cảnh pháp luật.
\end{enumerate}

\item Mức độ sai số trong Pre-FS cao hơn FS vì:
\begin{enumerate}[label=\Alph*.]
    \item Nhà đầu tư cố tình làm sai.
    \item Dữ liệu ở trạng thái tĩnh và mức độ chi tiết trung bình.
    \item Pre-FS không nghiên cứu thị trường.
    \item Pre-FS không cần thẩm định.
\end{enumerate}

\item Động từ phản ánh đúng nhất tính chất của "Thẩm định dự án" là:
\begin{enumerate}[label=\Alph*.]
    \item Mô tả. \item Đánh giá và Phản biện. \item Ghi chép. \item Phác thảo.
\end{enumerate}

\item Dự án thiếu phần "Nghiên cứu lợi ích KT-XH" trong FS sẽ gặp khó khăn gì?
\begin{enumerate}[label=\Alph*.]
    \item Không vay được vốn ngân hàng thương mại.
    \item Khó nhận được sự chấp thuận chủ trương đầu tư từ Nhà nước.
    \item Không tính được khấu hao.
    \item Không xác định được giá bán.
\end{enumerate}

\item Khi ý tưởng rất rõ ràng, nhà đầu tư có thể làm gì để tối ưu tiến độ?
\begin{enumerate}[label=\Alph*.]
    \item Bỏ qua cả Pre-FS và FS.
    \item Gộp bước nghiên cứu tiền khả thi vào nghiên cứu khả thi.
    \item Chờ thanh lý mới lập báo cáo.
    \item Chỉ nghiên cứu bối cảnh pháp luật.
\end{enumerate}

\item "Nghiên cứu về tổ chức nhân sự" trong FS nhằm mục đích:
\begin{enumerate}[label=\Alph*.]
    \item Để biết tên từng nhân viên.
    \item Xác định sơ đồ bộ máy, nhu cầu và chi phí lao động vận hành.
    \item Tính thuế TNCN cho giám đốc.
    \item Quảng bá thương hiệu.
\end{enumerate}

\item Chất lượng của công tác chuẩn bị ảnh hưởng trực tiếp đến:
\begin{enumerate}[label=\Alph*.]
    \item Chỉ ảnh hưởng đến giai đoạn ý tưởng.
    \item Toàn bộ quá trình quản lý thực hiện và nguồn lực đầu tư sau này.
    \item Chỉ ảnh hưởng đến thanh lý.
    \item Không ảnh hưởng hiệu quả tài chính.
\end{enumerate}

\item Tại sao "Nghiên cứu thị trường" luôn đứng đầu trong FS?
\begin{enumerate}[label=\Alph*.]
    \item Vì dễ làm nhất.
    \item Vì thị trường quyết định đầu ra (doanh thu), là gốc rễ của mọi tính toán.
    \item Vì không tốn chi phí.
    \item Vì thị trường luôn cố định.
\end{enumerate}

\item Hậu quả "động" khi bỏ qua "Bối cảnh pháp luật" trong FS bất động sản là:
\begin{enumerate}[label=\Alph*.]
    \item NPV tăng vọt.
    \item Dự án bị đình chỉ thi công do sai quy hoạch, gây lãng phí vốn.
    \item Máy móc hỏng nhanh hơn.
    \item Không thể thanh lý sau 50 năm.
\end{enumerate}

\item Đặc điểm "Chưa chi tiết" của Pre-FS phục vụ mục tiêu gì?
\begin{enumerate}[label=\Alph*.]
    \item Đánh lừa nhà đầu tư.
    \item Nhanh chóng có căn cứ sàng lọc và tiết kiệm chi phí trước khi nghiên cứu sâu.
    \item Thay thế thiết kế kỹ thuật.
    \item Tính thuế TNDN chính xác.
\end{enumerate}

\end{enumerate}

\section{Phần 3: Tính toán chỉ số hiệu quả tài chính NPV, IRR và Payback}

\begin{enumerate}[label=\textbf{Câu \arabic*:}, leftmargin=*, resume]

\item Năm cuối (năm 4): LNTT: 800; Khấu hao: 200; Thuế TNDN: 20\%; Thanh lý: 150. Tính $NCF_4$?
\begin{enumerate}[label=\Alph*.]
    \item 840. \item 990. \item 790. \item 1.150.
\end{enumerate}

\item Quên cộng "Giá trị thanh lý" vào năm cuối khi tính NPV dẫn đến:
\begin{enumerate}[label=\Alph*.]
    \item NPV bị thổi phồng.
    \item NPV bị đánh giá thấp hơn thực tế.
    \item IRR tăng cao bất thường.
    \item Thời gian hoàn vốn ngắn lại.
\end{enumerate}

\item Vốn đầu tư 2.000. NCF 3 năm: 800, 900, 1.000. Với $r=10\%$, NPV bằng bao nhiêu?
\begin{enumerate}[label=\Alph*.]
    \item 700. \item 222,39. \item 182,57. \item 2.222,39.
\end{enumerate}

\item Nếu dự án có NPV = 0 tại $r = 12\%$, và NPV = 100 tại $r = 10\%$. Chỉ số IRR là:
\begin{enumerate}[label=\Alph*.]
    \item 10\%. \item 11\%. \item 12\%. \item Không xác định.
\end{enumerate}

\item NCF hàng năm không đổi là 400, vốn đầu tư ban đầu 1.500. Tính $PP$?
\begin{enumerate}[label=\Alph*.]
    \item 3 năm. \item 3,75 năm. \item 4,2 năm. \item 5 năm.
\end{enumerate}

\item Khi lạm phát tăng làm $r$ tăng từ 10\% lên 15\%, NPV sẽ:
\begin{enumerate}[label=\Alph*.]
    \item Tăng. \item Giảm. \item Không đổi. \item Tiến về vô cực.
\end{enumerate}

\item Tại sao "Khấu hao" được cộng ngược lại vào LNST để tính NCF?
\begin{enumerate}[label=\Alph*.]
    \item Vì là khoản chi bằng tiền mặt.
    \item Vì là chi phí kế toán, không phải dòng tiền thực chi.
    \item Để làm IRR thấp xuống.
    \item Để trốn thuế.
\end{enumerate}

\item Dự án có NPV > 0 nhưng Payback Period vượt quá vòng đời dự án. Đánh giá?
\begin{enumerate}[label=\Alph*.]
    \item Dự án cực kỳ an toàn.
    \item Hiệu quả kinh tế nhưng rủi ro thanh khoản cao, không khả thi thực tế.
    \item Chắc chắn thành công.
    \item Chỉ số NPV bị tính sai.
\end{enumerate}

\item "Chi phí cơ hội" của vốn tự có được tính toán thông qua:
\begin{enumerate}[label=\Alph*.]
    \item Tiền phạt ngân hàng.
    \item Việc lựa chọn tỷ suất chiết khấu ($r$) phù hợp.
    \item Chi phí nguyên liệu.
    \item Tiền lương ban quản lý.
\end{enumerate}

\item Doanh thu: 1.000; Chi phí vận hành (không khấu hao): 400; Khấu hao: 100; Thuế: 20\%. NCF = ?
\begin{enumerate}[label=\Alph*.]
    \item 480. \item 500. \item 580. \item 600.
\end{enumerate}

\item Chỉ số nào phản ánh "khả năng chống chịu lãi suất" của dự án?
\begin{enumerate}[label=\Alph*.]
    \item NPV. \item IRR. \item Payback Period. \item Lợi nhuận thuần.
\end{enumerate}

\item Đầu tư 1.000, NCF năm 1: 600, năm 2: 600. $r=10\%$, tính NPV?
\begin{enumerate}[label=\Alph*.]
    \item 41,32. \item 200. \item -41,32. \item 100.
\end{enumerate}

\item Sai lầm khi không chiết khấu dòng tiền tương lai là vi phạm nguyên tắc:
\begin{enumerate}[label=\Alph*.]
    \item Nhanh gọn.
    \item Giá trị thời gian của tiền.
    \item Làm NPV âm.
    \item Thổi phồng rủi ro.
\end{enumerate}

\item Nếu $IRR < r$ (chi phí vốn), kết luận thẩm định là:
\begin{enumerate}[label=\Alph*.]
    \item NPV > 0.
    \item NPV < 0, dự án không đạt mức sinh lời kỳ vọng.
    \item Hoàn vốn nhanh.
    \item Không có chi phí cơ hội.
\end{enumerate}

\item Vốn đầu tư 500. NCF 3 năm: 100, 200, 300. Tính $PP$?
\begin{enumerate}[label=\Alph*.]
    \item 2 năm. \item 2,67 năm. \item 3 năm. \item 3,5 năm.
\end{enumerate}

\item Nhược điểm của IRR khi dòng tiền đổi dấu nhiều lần là gì?
\begin{enumerate}[label=\Alph*.]
    \item Không tính được NPV.
    \item Có thể cho ra nhiều kết quả (đa nghiệm).
    \item Payback bằng 0.
    \item Chi phí chuẩn bị tăng.
\end{enumerate}

\item Năm cuối: Doanh thu 500, Chi phí tiền mặt 200, Khấu hao 50, Thuế 20\%, Thanh lý 80. NCF = ?
\begin{enumerate}[label=\Alph*.]
    \item 330. \item 380. \item 280. \item 410.
\end{enumerate}

\item "NPV là phương tiện, không phải mục đích cuối cùng" nghĩa là:
\begin{enumerate}[label=\Alph*.]
    \item NPV không quan trọng.
    \item NPV cần xem xét cùng rủi ro và các yếu tố phi tài chính.
    \item NPV luôn đúng.
    \item NPV chỉ dùng cho chuẩn bị.
\end{enumerate}

\item Khi so sánh 2 dự án loại trừ nhau, dự án A có NPV cao hơn, dự án B có IRR cao hơn. Chọn?
\begin{enumerate}[label=\Alph*.]
    \item Chọn dự án B (IRR).
    \item Chọn dự án A (NPV) vì tối ưu giá trị tài sản ròng.
    \item Chọn vốn thấp hơn.
    \item Loại cả hai.
\end{enumerate}

\item Đầu tư 1 tỷ. NCF năm 1 là 1,2 tỷ. $r=10\%$. NPV là bao nhiêu?
\begin{enumerate}[label=\Alph*.]
    \item 90,9 triệu. \item 100 triệu. \item 200 triệu. \item 1,09 tỷ.
\end{enumerate}

\end{enumerate}

\section{Phần 4: Thực hiện đầu tư, Vận hành và Biến cố dự án}

\begin{enumerate}[label=\textbf{Câu \arabic*:}, leftmargin=*, resume]

\item Dự án đạt "Peak Capacity" chậm 2 năm so với FS. Chỉ số nào bị ảnh hưởng nhất?
\begin{enumerate}[label=\Alph*.]
    \item Tổng mức đầu tư.
    \item Thời gian hoàn vốn ($PP$) bị kéo dài.
    \item Chi phí lập Pre-FS.
    \item Giá trị thanh lý.
\end{enumerate}

\item "Thiết kế và lập dự toán thi công" thuộc giai đoạn nào?
\begin{enumerate}[label=\Alph*.]
    \item Chuẩn bị đầu tư.
    \item Thực hiện đầu tư.
    \item Vận hành.
    \item Thanh lý.
\end{enumerate}

\item Vốn đã huy động nhưng không thể triển khai do lỗi thủ tục gọi là:
\begin{enumerate}[label=\Alph*.]
    \item Vốn sinh lời cao.
    \item Vốn không sinh lời (Non-performing capital).
    \item Vốn lưu động.
    \item Vốn thanh lý.
\end{enumerate}

\item Công suất dự án "Giảm dần" ở cuối chu kỳ do:
\begin{enumerate}[label=\Alph*.]
    \item Chủ đầu tư chán.
    \item Hao mòn hữu hình và vô hình (lạc hậu công nghệ).
    \item Thuế tăng.
    \item NPV âm.
\end{enumerate}

\item Quản lý nguồn lực ở giai đoạn Thực hiện nhằm kiểm soát bao nhiêu \% vốn?
\begin{enumerate}[label=\Alph*.]
    \item 0,5 - 15\%. \item 50\%. \item 85 - 99,5\%. \item 100\%.
\end{enumerate}

\item "Tổ chức quản lý hoạt động kết quả đầu tư" quyết định:
\begin{enumerate}[label=\Alph*.]
    \item Chi phí lập Pre-FS.
    \item Hiệu quả và thời gian phát huy tác dụng của dự án.
    \item Số lượng ý tưởng mới.
    \item Giá trị vốn điều lệ.
\end{enumerate}

\item Khi thị trường giảm cầu đột ngột lúc đang vận hành, nhà quản lý nên:
\begin{enumerate}[label=\Alph*.]
    \item Ép máy chạy 100\%.
    \item Rà soát FS để điều chỉnh chiến lược hoặc thanh lý sớm.
    \item Cắt chi phí chuẩn bị.
    \item Tăng khấu hao.
\end{enumerate}

\item Giai đoạn "Sản xuất kinh doanh" đóng vai trò:
\begin{enumerate}[label=\Alph*.]
    \item Chi tiền mạnh nhất.
    \item Thu hồi vốn và tạo ra lợi nhuận.
    \item Lập dự toán.
    \item Sàng lọc ý tưởng.
\end{enumerate}

\item "Chạy thử và Nghiệm thu" là điểm giao thoa giữa:
\begin{enumerate}[label=\Alph*.]
    \item Ý tưởng và Chuẩn bị.
    \item Thực hiện đầu tư và Sản xuất kinh doanh.
    \item Vận hành và Thanh lý.
    \item Pre-FS và FS.
\end{enumerate}

\item Giá trị thanh lý thường bao gồm việc thu hồi "Vốn lưu động" vì:
\begin{enumerate}[label=\Alph*.]
    \item Là chi phí chìm.
    \item Tiền, tồn kho sẽ được thu hồi để hoàn trả chủ đầu tư khi dự án kết thúc.
    \item Nhà nước yêu cầu.
    \item Để tăng IRR ảo.
\end{enumerate}

\item Vòng đời dự án BOT cầu đường được xác định bởi:
\begin{enumerate}[label=\Alph*.]
    \item Khi cầu hỏng.
    \item Thời gian thu phí theo hợp đồng để hoàn vốn và có lời.
    \item Khi khởi công.
    \item Khi nộp Pre-FS.
\end{enumerate}

\item Lỗi "Thiên kiến chi phí chìm" xảy ra khi:
\begin{enumerate}[label=\Alph*.]
    \item Dừng dự án từ ý tưởng.
    \item Tiếp tục đổ tiền vào dự án kém hiệu quả vì "đã lỡ chi quá nhiều".
    \item Tính NPV quá kỹ.
    \item Thanh lý đúng hạn.
\end{enumerate}

\end{enumerate}
% --- PHẦN 1: ĐỀ THI --
\section{Câu hỏi tự luận}
\subsection{Xác định ý tưởng và phát hiện cơ hội đầu tư (2.5đ)}
\textbf{Nội dung:} Công ty A (khai khoáng) đang cân nhắc:
\begin{itemize}
    \item \textbf{PA 1:} Mở rộng dây chuyền truyền thống (nhân sự sẵn có, tài nguyên đã cấp phép nhưng áp lực thuế môi trường cao).
    \item \textbf{PA 2:} Sản xuất gạch không nung từ phế thải (ưu đãi thuế, lãi suất nhưng nguồn cung đầu vào phân tán).
\end{itemize}

\textbf{Yêu cầu:}
\begin{enumerate}[label=\alph*)]
    \item Dựa vào 06 tiêu chí xác lập ý tưởng, hãy đánh giá sự phù hợp của 2 phương án.
    \item Phân tích rủi ro "chuỗi cung ứng" đầu vào của PA 2 ảnh hưởng thế nào đến quyết định chọn ý tưởng?
\end{enumerate}

\subsection{Đánh đổi giữa chi phí chuẩn bị và độ chính xác (2.5đ)}
\textbf{Nội dung:} Dự án KDL Sinh thái 1.000 tỷ đồng.
\begin{itemize}
    \item Nhóm CEO: Chi 5 tỷ (0,5\%), làm trong 2 tháng.
    \item Nhóm Thẩm định: Chi 50 tỷ (5\%), làm trong 8 tháng (thuê tư vấn quốc tế).
\end{itemize}

\textbf{Yêu cầu:}
\begin{enumerate}[label=\alph*)]
    \item Tại sao chi 5\% ở giai đoạn chuẩn bị được coi là "khoản bảo hiểm" cho vốn thực hiện?
    \item Phân tích rủi ro "Vốn không sinh lời" và đưa ra phán quyết cuối cùng.
\end{enumerate}

\subsection{Phân tích động và quyết định đầu tư (2.5đ)}
\textbf{Nội dung:} Dự án C (linh kiện điện tử) gặp biến động khi lập báo cáo khả thi:
\begin{itemize}
    \item Đối thủ có công nghệ mới; Chi phí nhân sự tăng 25\%; Lãi suất vay tăng 2\%.
\end{itemize}

\textbf{Yêu cầu:}
\begin{enumerate}[label=\alph*)]
    \item Xác định các khối nghiên cứu (Modules) bị tác động trực tiếp.
    \item Tại sao phân tích ở trạng thái "Động" có độ tin cậy cao hơn để khẳng định cơ hội đầu tư?
\end{enumerate}

\subsection{Quản trị chu kỳ vận hành và thanh lý (2.5đ)}
\textbf{Nội dung:} Nhà máy D (vòng đời 8 năm). Năm 7-8 chi phí bảo trì chiếm 40\% doanh thu. Xuất hiện ý tưởng "Trung tâm thiết kế công nghệ cao".

\textbf{Yêu cầu:}
\begin{enumerate}[label=\alph*)]
    \item Xác định phân đoạn hiện tại của nhà máy D trong chu kỳ dự án.
    \item Ý nghĩa của việc nảy sinh ý tưởng mới đối với tính khép kín của chu kỳ dự án.
\end{enumerate}

\newpage
\section{BẢNG ĐÁP ÁN TỔNG HỢP}

\begin{longtable}{|c|c|m{3cm}|c|c|m{3cm}|}
\hline
\textbf{Câu} & \textbf{Đáp án} & \textbf{Cấp độ} & \textbf{Câu} & \textbf{Đáp án} & \textbf{Cấp độ} \\ \hline
1 & C & Phân tích & 33 & B & Vận dụng cao \\ \hline
2 & B & Vận dụng & 34 & B & Phân tích \\ \hline
3 & B & Vận dụng & 35 & B & Vận dụng cao \\ \hline
4 & B & Phân tích & 36 & C & Vận dụng \\ \hline
5 & B & Phân tích & 37 & B & Vận dụng \\ \hline
6 & B & Vận dụng & 38 & B & Phân tích \\ \hline
7 & B & Phân tích & 39 & B & Phân tích \\ \hline
8 & B & Phân tích & 40 & B & Phân tích \\ \hline
9 & B & Vận dụng & 41 & B & Phân tích \\ \hline
10 & B & Vận dụng & 42 & B & Vận dụng cao \\ \hline
11 & B & Phân tích & 43 & B & Phân tích \\ \hline
12 & B & Vận dụng & 44 & A & Vận dụng cao \\ \hline
13 & A & Vận dụng & 45 & B & Phân tích \\ \hline
14 & B & Phân tích & 46 & B & Phân tích \\ \hline
15 & B & Phân tích & 47 & B & Vận dụng cao \\ \hline
16 & B & Phân tích & 48 & B & Phân tích \\ \hline
17 & B & Vận dụng & 49 & A & Vận dụng cao \\ \hline
18 & C & Phân tích & 50 & B & Phân tích \\ \hline
19 & B & Phân tích & 51 & B & Phân tích \\ \hline
20 & B & Vận dụng & 52 & A & Vận dụng cao \\ \hline
21 & B & Phân tích & 53 & B & Phân tích \\ \hline
22 & B & Phân tích & 54 & B & Vận dụng \\ \hline
23 & B & Vận dụng & 55 & B & Phân tích \\ \hline
24 & B & Phân tích & 56 & B & Vận dụng \\ \hline
25 & B & Phân tích & 57 & C & Vận dụng \\ \hline
26 & B & Phân tích & 58 & B & Phân tích \\ \hline
27 & B & Vận dụng & 59 & B & Phân tích \\ \hline
28 & B & Vận dụng & 60 & B & Vận dụng \\ \hline
29 & B & Phân tích & 61 & B & Vận dụng \\ \hline
30 & B & Phân tích & 62 & B & Phân tích \\ \hline
31 & B & Phân tích & 63 & B & Phân tích \\ \hline
32 & B & Phân tích & 64 & B & Phân tích \\ \hline
\end{longtable}

\vspace{1cm}

\newpage

% --- PHẦN 2: ĐÁP ÁN VÀ GIẢI THÍCH CHI TIẾT ---
\begin{center}
    \large \textbf{ĐÁP ÁN VÀ HƯỚNG DẪN CHẤM CHI TIẾT}
\end{center}

\subsection{Tình huống 1: Chiến lược xác lập ý tưởng}
\begin{itemize}
    \item \textbf{Phân tích 06 tiêu chí:} Sinh viên cần nêu được: (1) Mục tiêu chủ thể, (2) Lợi thế cạnh tranh, (3) Sự phù hợp pháp luật, (4) Tính khả thi kỹ thuật, (5) Hiệu quả kinh tế, (6) Tác động môi trường.
    \item \textbf{So sánh:} PA 1 mạnh về (1, 2, 4) nhưng yếu về (3, 6). PA 2 mạnh về (3, 6) nhưng gặp rủi ro lớn ở tiêu chí (4) - Kỹ thuật/Cung ứng.
    \item \textbf{Trọng tâm:} Rủi ro nguồn cung phân tán là "điểm nghẽn" (bottleneck). Nếu không giải quyết được bằng hợp đồng bao tiêu, ý tưởng sẽ bị bác bỏ dù có ưu đãi thuế.
\end{itemize}

\subsection{Tình huống 2: Trade-off chi phí}
\begin{itemize}
    \item \textbf{Lập luận bảo vệ vốn:} Chi phí chuẩn bị ($C_p$) tỉ lệ nghịch với rủi ro sai lệch ($R$). Việc chi 5\% để giảm thiểu rủi ro cho 95\% vốn thực hiện ($V_{th}$) là tư duy quản trị rủi ro đúng đắn.
    \item \textbf{Vốn không sinh lời:} Giai đoạn này dòng tiền $CF < 0$. Tuy nhiên, nếu tiết kiệm chi phí chuẩn bị, dự án dễ rơi vào tình trạng "Chi phí chìm" (Sunk cost) khi phải sửa đổi thiết kế giữa chừng.
    \item \textbf{Phán quyết:} Chọn nhóm Thẩm định nhưng cần kèm KPI tiến độ để kiểm soát thời gian.
\end{itemize}

\subsection{Tình huống 3: Phân tích động và tài chính}
\begin{itemize}
    \item \textbf{Modules tác động:} Thị trường (Đối thủ), Tổ chức nhân sự (Chi phí chuyên gia), Tài chính (Lãi suất).
    \item \textbf{Công thức vận dụng:} Khi lãi suất ($r$) tăng, giá trị hiện tại thuần ($NPV$) giảm:
    \[ NPV = \sum_{t=1}^{n} \frac{CF_t}{(1+r)^t} - I_0 \]
    \item \textbf{Giải thích tính Động:} Trạng thái động phản ánh sự thay đổi của các biến số theo thời gian và thị trường. Điều này giúp kiểm tra tính bền vững của $IRR$ (suất sinh lời nội bộ). Nếu $IRR$ vẫn lớn hơn chi phí sử dụng vốn ($WACC$) sau biến động, dự án mới thực sự an toàn.
\end{itemize}

\subsection{Tình huống 4: Chu kỳ khép kín}
\begin{itemize}
    \item \textbf{Phân đoạn:} "Giai đoạn suy giảm và thanh lý". Do chi phí bảo trì quá cao (40\%), hiệu quả vận hành thực tế không còn.
    \item \textbf{Tính chu kỳ:} Việc hình thành ý tưởng mới ngay tại giai đoạn kết thúc dự án cũ tạo nên một hình xoắn ốc đi lên. Dự án sau kế thừa hạ tầng/thương hiệu của dự án trước nhưng ở trình độ công nghệ cao hơn, giúp doanh nghiệp tái sinh chu kỳ kinh doanh.
\end{itemize}

\rule{\textwidth}{1pt}

\begin{center}
\textit{--- Hết ---}
\end{center}

